\documentclass{article}
\usepackage[T1]{fontenc}
\usepackage{lmodern}
\usepackage{graphicx}
\usepackage{amsmath,amssymb}
\usepackage[a4paper,margin=2cm]{geometry}
\usepackage[absolute,overlay]{textpos}
\setlength{\TPHorizModule}{1mm}
\setlength{\TPVertModule}{1mm}
\usepackage[indonesian]{babel}
\usepackage{hyperref}
\setlength{\emergencystretch}{2em}
% unit: Halaman 10

\begin{document}

% === Halaman 10 (python/native) ===

10

• Contoh 1: 

 Plainteks: onetimepad

 Kunci:       tbfrgfarfm

 Misalkan A = 0, B = 1, …, Z = 25. 

 cipherteks: HOJKOREGHP

 yang dalam hal ini diperoleh sebagai berikut:


(‘o’ + ‘T’) mod 26 = H


(‘n’ + ‘B’) mod 26 = O


(‘e’ + ‘F’) mod 26 = J,  dst

\end{document}
